\documentclass[11pt]{article}
\usepackage{../../styles/activity}

\usepackage{xr}
\externaldocument{0-MR}

\lhead{}
%\chead{\textbf{\Large{\hspace{0pt}Beginning Activities for Section~3.3}}\\\prevhead}
\bahead{3.3}
\rhead{}
\lfoot{}
\rfoot{}
\cfoot{\hspace{0pt}\scalebox{0.4}{\includegraphics{cc-by-nc-sa.eps}}}

\begin{document}
\input table
\subsection*{Beginning Activity 1 (Proof by Contradiction)}
The main purpose of this beginning activity is to provide a logical justification for the method of proof know as ``proof by contradiction.''  A different justification for this method of proof is given in Exercise~1 of this section.
\begin{enumerate}
\item The truth table for $( {P \vee \mynot P} )$ shows that this statement is always true, and the truth table for $( {P \wedge \mynot P} )$ shows that this statement is always false.

\item
$$
\BeginTable
\BeginFormat
| c | c | c | c | c | c |
\EndFormat
\_6
| $P$  |  $Q$    \|6  $P \to Q$  |  $\mynot (P \to Q)$  |  $\mynot Q$ | $P \wedge \mynot Q$ | \\+22 \_6
          | T | T \|6 T | F | F | F | \\ 
         | T | F  \|6 F | T | T | T | \\ 
          | F | T \|6 T | F | F | F | \\ 
          | F | F \|6 T | F | T | F |  \\ \_6
\EndTable
$$

This shows that  $\mynot  ( {P \to Q} )$  is logically equivalent to  
$P \wedge \mynot  Q$.

\item If $x = 2$, then
\[
\frac{1}{x(1 - x)} = \frac{1}{2(-1)} = -\frac{1}{2},
\]
and $- \dfrac{1}{2} < 4$.  So $x = -2$ is a counterexample that shows the statement is false.

\item The negation of the statement is:
\begin{center}
There exists a real number $x$ such that $0 < x < 1$ and $\dfrac{1}{x(1 - x)} < 4$.
\end{center}
So for a proof by contradiction, we would assume that there exists a real number $x$ such that $0 < x < 1$ and $\dfrac{1}{x(1 - x)} < 4$.
\end{enumerate}
\hbreak

\newpage
\subsection*{Beginning Activity 2 (Constructing a Proof by Contradiction)}
\begin{enumerate}
  \item Since we have assumed that $x > 0$ and $y > 0$, we know that $xy > 0$.  So if we multiply both sides of the inequality $\dfrac{x}{y} + \dfrac{y}{x} \leq 2$, we obtain
\begin{align*}
\left( \frac{x}{y} + \frac{y}{x} \right) xy &\leq 2xy \\
\left( \frac{x}{y} \right) \cdot xy + \left( \frac{y}{x} \right) \cdot xy &\leq 2xy \\
x^2 + y^2 &\leq 2xy
\end{align*}
If we now subtract $2xy$ from both sides of the last inequality, and then factor the left side, we obtain
\begin{align*}
x^2 - 2xy + y^2 &\leq 0 \\
\left( x - y \right)^2 &\leq 0
\end{align*}

\item Another assumption that was made was that $x \ne y$.  This implies that $x - y \ne 0$ and hence, that 
$(x - y)^2 > 0$.  This contradicts the last inequality in part~(1).
\end{enumerate}
\hbreak

\end{document}

